\textbf{COLUNI 2024 -- Questão 30}\par\medskip Leia o texto a seguir.

A alimentação para corredores é o fator-chave para que o praticante, iniciante ou não, obtenha bons resultados e evite lesões. A corrida é um esporte que exige muita energia na sua prática. Por isso, é necessário suprir os estoques do organismo antes de sair para correr e repor o que foi perdido, logo após o treino.

Na alimentação pré-treino (ou pré-corrida, se preferir) é importante que o corredor procure por alimentos energéticos, como forma de ofertar combustível para o organismo suportar todo o percurso. A banana, a aveia, os cereais integrais, o pão integral e as raízes, como a batata doce ou aipim, são boas opções para serem ingeridas antes de encarar uma corrida.

Vale ressaltar que existem estratégias diferentes que podem ser adotadas por determinados períodos com o objetivo de potencializar e otimizar a utilização de gorduras durante o exercício e, assim, diminuir o consumo do glicogênio (reserva energética do músculo). Essas estratégias, geralmente, utilizam da redução de carboidratos e maior quantidade de gorduras boas antes do treino, como, por exemplo, abacate, oleaginosas, açaí (sem xarope) e até mesmo chocolate amargo.

Na alimentação pós-treino (ou pós-corrida, se preferir) o corredor deve pensar que precisa recuperar o que foi perdido e ainda dar nutrientes para que o organismo consiga reconstruir o que foi utilizado em demasia. Dessa forma, alimentos ricos em proteínas, como barras de proteínas, ovos cozidos, atum, sardinha e carne ou frango, são ótimas fontes alimentares pós-corrida. Para quem não tem alergia ou sensibilidade ao leite, por exemplo, pode utilizá-lo, bem como os seus derivados ou suplementos proteicos.

Disponível em: https://www.oceandrop.com.br/blog/6-dicas-para-corredores-iniciantes. Acesso em: 28 Ago. 2023. (Adaptado)

Com base no texto acima, responda às questões 27 a 30.

O movimento do corpo é um mecanismo muito importante no processo evolutivo das espécies animais, uma vez que permitiu, por exemplo, a busca de alimentos e a fuga de predadores, conferindo maior sobrevivência e reprodução para os indivíduos. O movimento surgiu no filo Cnidário e está presente tanto na fase de pólipo, quanto na fase de medusa.

\par

No que diz respeito ao movimento do corpo no filo Cnidário e aos mecanismos necessários para que ele aconteça, é INCORRETO afirmar que, esses organismos apresentam uma função:\begin{enumerate}
\item motora constituída de células contráteis que permitem a contração e o movimento.
\item circulatória constituída de vasos responsáveis por distribuir os nutrientes e disponibilizar energia para o
      movimento.
\item respiratória em que o gás oxigênio necessário para a respiração celular é retirado diretamente da água.
\item nervosa constituída de células nervosas e sensoriais, capaz de receber estímulos do ambiente e enviar
      resposta às células contráteis.
\end{enumerate}
